\cotefile{156-9}
\scanpage{18}
\mspage{14}
% archive sheet number: 17

des polygones et des figures, leurs $\delta\Phi$, $\partial\Phi$, ne sont pas changés, à un détail près, concernant \deleted{les} \inserted{définition et terminologie concernant les} ``\uline{sommets} \uline{intérieurs} \deleted{de $\Phi$} ou ``internes'' de $\Phi$, propres ou impropres. En plus des sommets internes pris précédemment (\deleted{des} $s\in\delta\Phi$ t.q. $\operatorname{ord}(s,\Phi)=2$), il faut rajouter les \deleted{$s\in\delta\Phi$} $s\in\partial I$ tels que $\operatorname{ord}(s,\unsure{\Phi|\mathbb{\Phi}})=1$, \inserted{i.e. $s\in\partial\Phi$.} (Car si $s\in\unsure{\delta I|\partial I}$, on a tjrs $\operatorname{ord}(s,\Phi)\in\{0,1\}$, à l'exclusion de \deleted{\unsure{2}} $2$, \deleted{et cela se} \unsure{le disant} pour $\forall\,\Phi$ ou pour $\forall\,\Phi$ de la forme $\{\deleted{I}\unsure{X|\alpha|\,}\}$, où $\alpha\in\unsure{M_1|\mathcal M_1}$, \unsure{cela caractérise} les pts de $\delta I$). On dit qu'un tel sommet est \deleted{\unsure{lac|non}} \uline{redondant} s'il est propre, ou encore si (comme pour les sommets redondants proprement dits) il n'est pas dans $\delta\Phi^*$; dans le cas inverse, il est dit sommet lacunaire. Ceci dit, on a
\[
\Phi^*_{\mathrm{is}}=\Phi_{\mathrm{lac}}\qquad \Phi^*_{\mathrm{lac}}=\Phi_{\mathrm{is}}
\]
\[
\Phi_{\mathrm{red}}=\delta\Phi-\delta\Phi^*
\]
\[
\delta\Phi=\underbrace{\Phi_{\mathrm{red}}\cup\Phi_{\mathrm{lac}}}_{\Phi_{\mathrm{int}}}\cup\Phi_{\mathrm{is}}\cup
\underbrace{\partial\Phi}_{
 \underbrace{\partial_0\Phi}_{\substack{\Phi_0\cap\delta\Phi\\\text{sommets-bord}\\\text{propres}}}
 \amalg
 \underbrace{\partial'_0\Phi}_{\text{somm.b. impropres}}}
\]
\margin{Écrire réunion \uline{disjointe}
$\delta\Phi=\delta\Phi=\Phi_{\deleted{\unsure{\mathrm{red}}}\inserted{\mathrm{red}}}\cup\Phi_{\deleted{\unsure{\mathrm{red}}}\mathrm{is}}\cup\Phi_{\mathrm{lac}}\cup(\unsure{\partial^*|\delta^*}\Phi\deleted{\unsure{\delta\Phi\cap}}$
où $\partial^*\Phi=\partial\Phi\setminus\delta\Phi\cap\partial I$ (sommets-bords \uline{non} \uline{extérieurs})}
réunions disjointes, sauf pour \unsure{le dernier terme|les derniers termes} $\partial\Phi$ et les deux premiers, formant $\Phi_{\mathrm{int}}$, car
\[
\partial\Phi\cap\Phi_{\mathrm{int}}=\partial I\cap\deleted{\unsure{\delta I\cap}}(\deleted{\unsure{\delta\Phi}}\Phi_0\setminus\Phi_{\mathrm{is}})
\]
les points de $\delta\Phi$ qui sont intérieurs sont les \deleted{points} qui sont dans $\partial I$, sans être des points isolés de $\Phi$.
\scanpage{19}
\mspage{15}
% archive sheet number: 18

Du coup, l'intérêt de l'introduction de $\partial\Phi$ \deleted{des} (des sommets-bord de $\Phi$) devient plus problématique, d'autant plus que \deleted{ces sommets} \inserted{l'on n'a} \deleted{\unsure{bord} de $\Phi$ \unsure{ni}} \inserted{plus} évidemment $\partial\Phi=\partial\Phi^*$. \deleted{\unsure{Ci-vant|Ci-contre|à voir}}
\deleted{En fait, on a}
\[
\deleted{\partial\Phi^*\cap\partial\Phi=\partial\Phi\setminus\partial\Phi\cap\partial I}
\]
\deleted{et}
\[
\deleted{\partial\Phi=(\partial\Phi^*\cap\partial\Phi)\amalg\underbrace{\partial\Phi\cap\partial I}_{\substack{\Vert\\\Phi_{\mathrm{int}}\cap\partial I}}}
\]
\[
\deleted{\partial\Phi^*=(\partial\Phi^*\cap\partial\Phi)\amalg\Phi_{\mathrm{is}}\cap\partial I}
\]

On a
\[
\left.\begin{aligned}
\Phi_{\mathrm{is}}&=\Phi^*_{\mathrm{lac}}\\
\Phi_{\mathrm{lac}}&=\Phi^*_{\mathrm{is}}
\end{aligned}\right\}\quad\begin{gathered}\text{\deleted{\unsure{sommets}} lacunaires et s. isolés de $\Phi$ et de $\Phi^*$}\\\text{s'échangent}\end{gathered}
\]
\[
\begin{aligned}
\Phi_{\mathrm{red}}&=\delta\Phi\setminus\delta\Phi\cap\deleted{\unsure{\partial}}\delta\Phi^*\qquad \text{NB. }\Phi^*_{\mathrm{red}}=\partial I\setminus(\partial I\cap\partial\Phi)-\\
&=\bigl(\underbrace{\delta\Phi\setminus\delta\Phi^{*\deleted{*}}}_{\text{sommets évanescents}}\bigr)\amalg\underbrace{(\partial I\cap\partial\Phi)}_{\text{sommets redond}}
\end{aligned}
\]
\[
\deleted{\unsure{\Phi}}\unsure{\partial^*|\delta^*}\Phi=\partial^*\Phi^*\qquad\text{où}\quad\partial^*\Phi=\partial\Phi\setminus\delta\Phi\cap\partial I
\]

Ainsi dans la décomposition canonique
\[
\delta\Phi=\Phi_{\mathrm{red}}\cup\Phi_{\mathrm{is}}\cup\Phi_{\mathrm{lac}}\cup\partial^*\Phi
\]
les $2^{\mathrm{e}}$ et $3^{\mathrm{e}}$ termes s'échangent par dualité, le troisième \unsure{ne varietur|invariant|ne varie pas} par dualité, le premier disparaît par dualité -- mais \deleted{\unsure{réappar}} la partie \deleted{$\partial\Phi\cap\partial I$} $\Phi_{\mathrm{red}}\cap\partial I$ \deleted{\unsure{réapp}} de $\Phi_{\mathrm{red}}$ réapparaît par bidualité: \deleted{\unsure{sommets redon.}} Du coup le terme ``redondant'' pour ces sommets n'est pas tellement adéquat, \deleted{\unsure{on peut les appeler}} si on adopte la terminologie
\[
\text{sommets évanescents}\qquad \Phi_{\text{év}}\qquad\text{old somm.réd. }\Phi_{\mathrm{red}}
\]
on pourrait appeler ces sommets semi-évanescents. On trouve alors une décomposition
\scanpage{20}
\mspage{16}
% archive sheet number: 19

disjointe autoduale de
\[
\widehat\delta\Phi\overset{\mathrm{def}}=\delta\Phi\cup\partial I=\delta\Phi\cup\underbrace{(\partial I\setminus\unsure{\delta I|\partial I}\cap\delta\Phi)}_{=\Phi^*_{\text{sév}}}
\]
sév = semi-évanescents, ainsi.
\diagram{\begin{tikzcd}[ampersand replacement=\&,column sep=small,row sep=large]
\widehat\delta\Phi=\Phi_{\mathrm{is}} \arrow[rr,leftrightarrow,bend right=40] \arrow[rr,leftrightarrow,bend right=44] \arrow[rr,bend right=48] \& {}\cup{} \& \Phi_{\mathrm{lac}} \& {}\cup\Phi_{\mathrm{ev}} \arrow[d] \& {}\cup\Phi_{\mathrm{sev}} \arrow[r,leftrightarrow,bend right=30] \arrow[d,equal] \& {}\cup\Phi^{\unsure{*|'}}_{\mathrm{sev}} \arrow[d,equal] \& {}\cup\partial^*\Phi \arrow[loop below,leftrightarrow]\!\\
\&\&\&\substack{\text{disparaît}\\\text{par dualité}}\&\deleted{\unsure{\partial I\setminus\partial I\cap\delta\Phi^*}}\arrow[d,equal]\&\partial I\setminus\partial I\cap\delta\Phi\&\\
\&\&\&\&\deleted{\unsure{\delta\Phi\setminus\delta\Phi^*\cap\delta\Phi}}\&\&
\end{tikzcd}}
où
\[
\begin{aligned}
\Phi_{\mathrm{is}}&=\Phi^*_{\mathrm{lac}}\\
\Phi_{\mathrm{lac}}&=\Phi^*_{\mathrm{is}}\\
\partial^*\Phi_{\deleted{\unsure{\mathrm{sev}}}}&=\partial^*\Phi^*\\[1ex]
\Phi^*_{\mathrm{ev}}&=\Phi^{**}_{\mathrm{ev}}=\varnothing
\end{aligned}
\]

Pour les $s\in\partial I$, il y a quatre positions possibles p.r. à $\Phi$:
\[
\begin{aligned}
\text{a)}\quad s\in\Phi_{\mathrm{is}}&\Longleftrightarrow s\in\Phi^*_{\mathrm{lac}}\\
\text{b)}\quad s\in\Phi_{\mathrm{sev}}&\Longleftrightarrow s\notin\delta\Phi^*\\
\text{a$^*$)}\quad s\in\Phi_{\mathrm{lac}}&\Longleftrightarrow s\in\Phi^*_{\mathrm{is}}\\
\text{b$^*$)}\quad s\notin\delta\Phi&\Longleftrightarrow s\in\Phi^*_{\mathrm{sev}}
\end{aligned}
\]
\[
\begin{aligned}
\Phi_{\text{sév}}&=\{s\in\partial I\mid s\notin\delta\Phi^*\}.\\
&=\partial I\setminus\partial I\cap\delta\Phi^*\\
\Phi^*_{\mathrm{sev}}&=\{s\in\partial I\mid s\notin\delta\Phi\}\\
&=\partial I-\partial I\cap\delta\Phi
\end{aligned}
\]
% slanted
\margin{\uline{Cor 1} \deleted{\unsure{Omb}} $\unsure{\operatorname{Supp}^{\circ}|\operatorname{S-m}^{\circ}}(\Phi)=\operatorname{Supp}^{\circ}(\operatorname{Omb}^{\circ}\Phi)=$\par
\uline{Cor. 2} $\operatorname{Supp}^{\circ}\Phi=\operatorname{Omb}^{\circ}\Phi^{**}$}

On trouve alors
\[
\boxed{\unsure{\operatorname{cosupp}^{\circ}|\operatorname{cos-m}^{\circ}}(\Phi)=\operatorname{cosupp}^{\circ}(\operatorname{Omb}^{\circ}(\Phi))=\operatorname{Omb}^{\circ}(\Phi^*)}
\]

On peut appeler \uline{parties constructibles de $\unsure{M|\mathcal M}$} les parties de la forme $\operatorname{Omb}^{\circ}(\Phi)$, $\Phi$ polygone. \deleted{Il y a} L'application
\[
\Phi\longmapsto\operatorname{Omb}^{\circ}(\Phi)\qquad\operatorname{Polyg}(M)\longrightarrow\operatorname{Cons}(M)
\]
est bijective. En effet, on récupère $\Phi$
\scanpage{21}
\mspage{17}
% archive sheet number: 20

à partir de $\operatorname{Omb}^{\circ}(\Phi)=\Sigma$ par
\[
\left\{\begin{aligned}
\Phi_1&=\text{ensemble des éléments de }\deleted{\operatorname{Omb}^{\circ}(\Phi)}\Sigma\ \deleted{\unsure{tels}},\\
&\qquad\text{maximaux pour }\overset{\circ}{\ll},\text{ qui sont dans }\unsure{M_1|\mathcal M_1},\\
&\deleted{\Phi_0=\{s\in\unsure{\Sigma|M_0}\mid}\\
\Phi_{\mathrm{is}}&=\text{ensemble des éléments de }\Sigma,\text{ maximaux}\\
&\qquad\text{pour }\overset{\circ}{\ll},\text{ qui sont dans }M_0.\\
\Phi_0\setminus\Phi_{\mathrm{is}}&=\left(\bigcup_{X\in\Phi_1}\partial X\right)\cap\Sigma
 =\deleted{\{s\in\Sigma\cap M_0\}},\\
&=\{s\in\Sigma\cap M_0\mid\exists\,X\in\Phi_1\text{ t.q. }s\unsure{\triangleleft|\in\partial}X\}
\end{aligned}\right.
\]

On trouve que le cosupport d'une partie constructible de $M$ est encore constructible, et le passage au cosupport, \inserted{correspond} du côté des polygones, au passage à \deleted{\unsure{support}} le polygone complémentaire. Le passage au \inserted{support correspond \unsure{plutôt} au passage à l'\uline{enveloppe} (\uline{\unsure{irredondante|irrédundante}}) de la \uline{figure envisagée}.}
\margin{$\unsure{\flat|b}\Phi^{**}$}

Considérons un drapeau $T\in\operatorname{Drap}(I)$, $T\supset\partial I$,
\[
t_0=\omega_0<t_1\ldots<t_n=\omega_1
\]
($n=\text{longueur du drapeau}$), les strates constructibles sont
\[
\underset{\substack{\Vert\\\omega_0}}{\{t_0\}},\ \deleted{\unsure{\partial}}I_{t_0t_1},\ t_1,\ I_{t_1,t_2}\quad\cdots\quad I_{t_{\unsure{n-1}},t_n},\ \underset{\substack{\Vert\\\omega_1}}{t_n}.
\]
On appelle polygone \uline{subordonné} au drapeau $T$, toute partie \inserted{$\Phi$} de cet ensemble $\Phi_T$. On trouve alors (sans se préoccuper du support) ceci, $\forall\,\Phi\subset\Phi_T$, on désigne par $C\Phi$ son complémentaire dans $\Phi_T$
\scanpage{22}
\mspage{18}
% archive sheet number: 21

\[
\underbrace{(C\Phi)^{**}}_{\substack{\text{figure}\\\text{\unsure{irredondante|irrédundante}}\\\text{associée à }C\Phi}}=\Phi^*
\]
d'où, \unsure{en vertu} de la formule fond. plus \inserted{haut}
\[
\boxed{\begin{aligned}
\unsure{\operatorname{Cosupp}^{\circ}|\operatorname{Eosupp}^{\circ}}\Phi&=\unsure{\operatorname{cosupp}^{\circ}|\operatorname{supp}^{\circ}}\operatorname{Omb}^{\circ}\Phi=\operatorname{supp}^{\circ}C\Phi.\\
&=\deleted{\operatorname{cosupp}^{\circ}\unsure{\operatorname{supp}}}
\end{aligned}}
\]
Soit alors, pour $\forall\,\Phi\subset\Phi_T$,
\[
\unsure{\mathbb{S}|S|\Sigma}_{\Phi}\subset\unsure{M|\mathcal M},\quad\mathbb{S}_{\Phi}=\operatorname{supp}^{\circ}(\unsure{\Phi|\mathbb{\Phi}})\unsure{=|\deleted{=}}\operatorname{supp}^{\circ}(\operatorname{Omb}^{\circ}\Phi\unsure{)|\,}
\]
On trouve une application
\[
\Phi\longmapsto S_{\Phi}\qquad\unsure{\mathscr P|\mathcal P}(\Phi_T)\longrightarrow\underset{\substack{\text{ens. des}\\\text{supports}\\\text{constructibles}\\\text{de }M}}{\Sigma_{\mathrm{cons}}(M)}
\]
Cette application transforme \unsure{tautologiquement} sommes (i.e. $\operatorname{Sup}$ pour $\subset$) en $\operatorname{Sup}$ dans $\Sigma_{\mathrm{cons}}$. La formule précédente signifie qu'elle commute à $C$. Il en résulte qu'elle commute aux intersections. \uline{C'est l'axiome des supports At supp 1}.

Deux \deleted{parties} polygones \inserted{$\Phi,\Psi$} sont dits \uline{commensurables} si $\delta\Phi\cup\delta\Psi$ est un drapeau. Cela \deleted{signifie} implique que leurs
\scanpage{23}
\mspage{19}
% archive sheet number: 22

supports \inserted{$\operatorname{supp}^{\circ}$} sont ``\uline{subordonnés}'' à un même drapeau $\delta\Phi\cup\delta\Psi\cup\partial I=T$. [Idem pour un ensemble quelconque de polygones, ou de parties constructibles de $\unsure{M|\mathcal M}$ (p.ex. des supports)]\deleted{\unsure{ceci}} -- quitte à subdiviser les intervalles polygonaux dans $\Phi,\Psi$ (ce qui ne change pas le $\operatorname{supp}^{\circ}$).

Toutes les opérations \inserted{usuelles} qu'on peut faire, sur une famille quelconque d'\deleted{$\operatorname{supp}^{\circ}$}éléments de $\Sigma_{\mathrm{cons}}$ mutuellement commensurables, se ramènent aux opérations $\cup,\cap,C$ dans un $\unsure{\mathscr P|\mathcal P}(\Phi_T)$ convenable.
(NB si $T$ de long. $n$, avec $n$ intervalles et $n+1$ sommets ($n\geq1$)
\[
\operatorname{card}\Phi_T=2n+1
\]
et
\[
\underbrace{\operatorname{Card}\mathscr P(\Phi_T)}_{\substack{\text{nb de supports}\\\text{subordonnés à }T}}=2^{2n+1}.
\]

\textcircled{2} Cas d'un $I$ où on ne se donne pas $\omega_0,\omega_1$. On doit le voir comme déduit du cas précédent $(I,\partial I)$, en \unsure{passant au}
\scanpage{24}
\mspage{20}
% archive sheet number: 23

``\unsure{magasin} résiduel'' ou ``complémentaire'' p.r. à une des extrémités de $I$, ou des deux. Supposons p.ex. \deleted{qu'il n'y ait} qu'on regarde le $(I,\partial I)\setminus\{\omega_0\}$ (façon de parler) -- ou, plus précisément, le sous-magasin
\[
\deleted{\unsure{M_{\overset{\circ}{\ll}}}}\operatorname{Omb}^{\circ}(\{I_{\omega_0,\omega_1},\omega_1\})=\unsure{M|\mathcal M}\setminus\omega_0
\]
(partie de $M$ fermée pour $\unsure{\omega_0|\overset{\circ}{\ll}}$), on désire l'interpréter en termes de
\[
(\overbrace{I\setminus\{\omega_0\}}^J,\omega_1)
\]
un ordonné $J$, \uline{muni} d'un plus gd élément $\omega_1$. On garde toutes les strates de $M(I,\partial I)$, à l'exception de $\omega_0$. Mais les strates de la forme \inserted{avec $\omega_0<a$ i.e. $a\in J$} $I_{\omega_0,a}$, \deleted{s'écrivent} maintenant
\[
J_{<a}
\]
\margin{On ne suppose plus ici que $a$ ne soit pas \unsure{minimal}, cf plus bas}
Les relations $\unsure{\trianglelefteq|\triangleleft},\leq,\overset{\circ}{\ll},\mathrel{|\circ|}$ \deleted{et} dans l'ens.
\[
M(J,\omega_1)\deleted{\unsure{=}}\quad(=M(I,\partial I\setminus\{\omega_0\})
\]
pourront alors s'expliciter directement, en termes de $J,\omega_1$, sans réf. \inserted{explicite} à $(I,\partial I)$, en travaillant avec les strates
\[
a,\ J_{<a},\ \unsure{\mathbb I|I|J}_{a,b}\qquad a<b,\quad a,b\in J.
\]
Si p.ex. $I=\{\omega_0,\omega_1\}$, on aura $J=\{\omega_1\}$,
\scanpage{25}
\mspage{21}
% archive sheet number: 24

et
\[
\unsure{M|\mathcal M}(J,\omega_1)=\{\omega_1,J_{<\omega_1}\}
\]
magasin qui consiste en \uline{deux} éléments seulement, et \uline{c'est une moquette}. Ici, on a
\[
|J_{<\omega_1}|^{\circ}=\varnothing,\quad\text{mais bien sûr}\quad J^{\circ}_{<\omega_1}\neq\varnothing.
\]

\deleted{Choix dual} Traitement dual pour
\[
M(I,\partial I)\setminus\{\omega_1\}\overset{\mathrm{def}}=M(J,\omega_0)
\]
où $\omega_0$ est donné en \uline{tant que} plus petit élément de $J\equiv I\setminus\{\omega_1\unsure{)|\}}$. Cette fois, c'est les éléments
\[
J_{>a}\qquad(a\in J)
\]
qu'il faut introduire, en plus des ``lieux'' $a\in J$ et des ``intervalles finis'' \unsure{de $J$}
\[
J_{a,b}.
\]

Prenons enfin le cas du magasin résiduel de $\partial I$ dans $M$
\[
M(I,\partial I)\setminus\partial I\overset{\mathrm{def}}=M(J)\qquad J=I\setminus\partial I
\]
Cette fois, il faut introduire à la fois les
\[
J_{<a}\quad\text{et les}\quad J_{>a}\qquad(a\in J),
\]
% slanted
\margin{\unsure{ceci}, y compris si $a$ est minimal ou maximal dans $J$!}
pour les ``traces'' des intervalles \deleted{\unsure{en}} $I_{\omega_0,a},I_{a,\omega_1}$, et \deleted{de plus} enfin, de plus, un objet noté
\[
[\unsure{\mathbb J|J}]
\]
qui tient lieu de $I_{\omega_0,\omega_1}$. Les \unsure{introductions}
